% =====================================================================
% Section 5.2 "Design ablations" -- MEASURED, 15 November 2026.
% Replaces the \todo{run} table and the \todo paragraph in
% SC4001_report_restructured (1).tex, lines ~823-872.
%
% >>> READ THIS BEFORE PASTING -------------------------------------
% The \todo paragraph in the draft predicts four outcomes. THREE OF THEM
% ARE CONTRADICTED by the runs. Do not paste the predicted narrative.
%
%   predicted: "the classic overfitting curve" against hidden width
%   measured : a PLATEAU. linear -0.053, 16 units -0.008, 32-16 -0.004,
%              128-64 -0.004. Quadrupling the width changes nothing.
%              There is no overfitting-with-width curve to draw.
%
%   predicted: "widening the feature set buys variance, not signal"
%   measured : true for AF3 features (5 -> 21 moves -0.029 -> -0.067),
%              FALSE overall -- the 16 design-side features recover it
%              to -0.004. The story is which features, not how many.
%
%   predicted: "the pairwise loss was validated against the pointwise
%              alternative"
%   measured : they are INDISTINGUISHABLE, and pointwise BCE is
%              nominally ahead (+0.005 vs -0.004). You may justify the
%              pairwise loss by task match; you may NOT say the
%              comparison validated it.
%
%   predicted: within-target normalisation matters
%   measured : CONFIRMED, and it is the only factor whose interval
%              excludes zero. Pooling costs 0.126 AP.
%
% All runs: features_folds.csv, 3,532 designs, 13 receptors, LOTO,
% 5 model seeds averaged per fold, 2,000-resample paired bootstrap.
% Reproduce with reranker/run_ablations.sh.
% =====================================================================

\subsection{Design ablations: why the learned head cannot help}\label{sec:r-ablations}
A null result is only informative if the model was given a fair chance. We therefore ablate the
four design choices of \cref{sec:mlp} --- capacity, feature count, loss, and input normalisation
--- holding the leave-one-target-out protocol and the feature table fixed and changing one
choice at a time (\cref{tab:designablation}). Every run is CPU-only and takes about 80\,s.

\begin{table}[H]
  \centering
  \caption{Re-ranker design ablations under the same leave-one-target-out protocol. Each row
  changes one choice from the reference configuration. Intervals are target-level paired
  bootstraps against \ipsaemin{} with 2,000 resamples; $^{*}$ marks an interval excluding zero.}
  \label{tab:designablation}
  \small
  \begin{tabular}{llcc}
    \toprule
    Factor & Setting & Mean per-target AP $\uparrow$ & \dAP{} vs.\ \ipsaemin{} [95\% CI]\\
    \midrule
    --- & reference ($37\rightarrow32\rightarrow16$) & 0.545 & $-0.004$ [$-0.071$, $+0.074$]\\
    \midrule
    \multirow{3}{*}{Capacity}
      & linear ($37\rightarrow1$)               & 0.496 & $-0.053$ [$-0.132$, $+0.016$]\\
      & one hidden layer ($37\rightarrow16$)    & 0.540 & $-0.008$ [$-0.078$, $+0.058$]\\
      & wider ($37\rightarrow128\rightarrow64$) & 0.545 & $-0.004$ [$-0.074$, $+0.076$]\\
    \midrule
    \multirow{3}{*}{Feature count}
      & 5 AF3 features     & 0.519 & $-0.029$ [$-0.071$, $+0.009$]\\
      & 21 AF3 features    & 0.481 & $-0.067$ [$-0.144$, $-0.009$]$^{*}$\\
      & 37 (AF3 $+$ input) & 0.545 & $-0.004$ [$-0.071$, $+0.074$]\\
    \midrule
    \multirow{2}{*}{Loss}
      & pairwise logistic (RankNet) & 0.545 & $-0.004$ [$-0.071$, $+0.074$]\\
      & pointwise BCE               & 0.554 & $+0.005$ [$-0.038$, $+0.053$]\\
    \midrule
    \multirow{2}{*}{Normalisation}
      & $z$-score within target & 0.545 & $-0.004$ [$-0.071$, $+0.074$]\\
      & $z$-score pooled        & 0.418 & $-0.130$ [$-0.240$, $-0.039$]$^{*}$\\
    \bottomrule
  \end{tabular}
\end{table}

\paragraph{Only one design choice matters, and it is not the architecture.} Pooling the
$z$-scores across targets instead of within them costs 0.126 AP and is the single largest effect
in the table --- the only architectural factor whose interval excludes zero. Pooling leaves the
per-target score offsets in the data, so the model can lower its training loss by learning which
receptor a design belongs to; at test time the held-out receptor is by construction one it has
never seen, and that capacity is worse than wasted. This is the grouped-data failure of
\cref{sec:r-ladder} reappearing inside our own model.

\paragraph{Capacity does not bind.} A linear probe is the weakest arm ($-0.053$), one hidden
layer of 16 units already recovers almost all of the gap ($-0.008$), and quadrupling the hidden
width to $128\rightarrow64$ changes the result not at all ($-0.004$, identical to the reference
to three decimals). We expected the classic overfitting curve on 13 effective training units and
did not find one: the curve is a plateau. The model is not capacity-limited, and it is not
capacity that stops it beating \ipsaemin{}.

\paragraph{It is which features, not how many.} Within the AF3 block, adding features makes the
model worse --- five confidence scores give $-0.029$, the full 21 give $-0.067$, the only arm
besides pooled normalisation whose interval excludes zero. The 16 design-side features then
recover the loss ($-0.004$), and they are the only features computed \emph{before} any
prediction. So the widely-stated reading, that more features buy variance rather than signal, is
only half right here: more \emph{AF3} features buy variance, because they are correlated
summaries of one confidence head, while features from a different source add genuine
information. With 13 receptors as the unit of replication the effective sample size for
generalisation is 13, not 3,532, and that budget is spent quickly on redundant inputs.

\paragraph{The loss function is not the lever either.} Pointwise BCE and the pairwise RankNet
loss are statistically indistinguishable ($+0.005$ against $-0.004$; the intervals overlap over
almost their whole length). We retain the pairwise loss because it matches the task --- ranking
within a target, not calibrated classification across them --- but we report plainly that the
comparison does not establish it as superior. The BCE arm is nominally the best configuration we
ran, and its interval still contains zero, so it does not beat \ipsaemin{} either.

\paragraph{What the grid shows.} No configuration we tried beats the score it re-ranks. The
factors that could be blamed on us --- too little capacity, the wrong loss, too few features ---
all move the result by less than the width of its own confidence interval, while the one factor
that does move it is a property of how the data is grouped. That is the evidence for
\cref{sec:r-mlp}'s reading: the ceiling is the information AF3's confidence head exposes and the
13 receptors available to generalise over, not the model placed on top.
